\section{Discussion and Limitations}

\subsection{Model 1 Interpretation}

The experiments show that a relatively small scratch CNN can already achieve
strong results on EuroSAT RGB. The architecture alone was not the only important
part. The training configuration also mattered. The best validation result was
obtained only after combining augmentation, dropout, weight decay, and SGD with
momentum.

The binary classification task
reached a test macro-F1 of 0.9795, while the 10-class classification reached
0.9519. This is expected because the binary task groups several similar classes
together. Vegetation classes such as Pasture, PermanentCrop,
HerbaceousVegetation, and AnnualCrop can be difficult to separate in RGB images.

One advantage of the chosen procedure is the clean separation between training,
validation, and test data. Normalization was computed only from the training
data, and the test split was not used for hyperparameter selection. Therefore,
the final test scores are more meaningful than they would be if the test split
had been used repeatedly during model tuning.

\subsection{ResNet Transfer Learning Interpretation}

The ResNet-50 transfer learning model achieves a test macro-F1 of 0.9812 on the
10-class task, compared to 0.9519 for the scratch CNN. The improvement is
moderate in absolute terms, but the data efficiency experiments reveal a much
larger story.

The frozen ResNet baseline, which trains only the linear classification head
(20{,}490 parameters) while keeping all backbone weights fixed, reaches a test
accuracy of 93.7\%. This means that ImageNet features alone, without any
fine-tuning, already outperform the fully trained scratch CNN at 500 and 2{,}000
training images. It confirms that the backbone has learned feature representations
that transfer directly to satellite land-use classification.

Fine-tuning \texttt{layer4} on top of the frozen features adds a further 4.4
percentage points at full data, pushing test accuracy from 93.7\% to 98.2\%.
This marginal cost (adapting roughly 14.7 million parameters instead of 20{,}490)
is justified by the consistent gain across all data sizes.

One limitation of the transfer learning approach is the increased computational
cost per epoch. With 224$\times$224 inputs and a 50-layer backbone, each epoch
takes approximately 151 seconds on Apple Silicon MPS, compared to 8 seconds for
the scratch CNN. However, the ResNet requires far fewer epochs to converge
(15 vs.\ 42 for the scratch CNN at full data), so the total wall-clock training
time is roughly 37 minutes versus 6 minutes. For applications where training
time matters more than final accuracy, the scratch CNN remains a practical choice.

\subsection{Model Comparison}

Table \ref{tab:model_comparison} compares the scratch CNN with the two Model 2
transfer-learning backbones across both tasks. Both pretrained backbones
outperform the scratch CNN on both tasks: ResNet-50 achieves the best 10-class
result, while EfficientNet-B0 is marginally the best on the binary task.

\begin{table}[H]
\centering
\caption{Final model comparison on the full test split ($n=4050$).}
\label{tab:model_comparison}
\small
\begin{tabular}{llrrrr}
\toprule
Model & Task & Val Acc. & Val F1 & Test Acc. & Test F1 \\
\midrule
Scratch CNN        & 10-class & 0.9573 & 0.9557 & 0.9531 & 0.9519 \\
Scratch CNN        & Binary   & 0.9795 & 0.9795 & 0.9795 & 0.9795 \\
EfficientNet-B0    & 10-class & 0.9793 & 0.9784 & 0.9726 & 0.9719 \\
EfficientNet-B0    & Binary   & 0.9961 & 0.9961 & 0.9941 & 0.9941 \\
ResNet-50 (layer4) & 10-class & 0.9780 & 0.9770 & 0.9817 & 0.9812 \\
ResNet-50 (layer4) & Binary   & 0.9943 & 0.9943 & 0.9938 & 0.9938 \\
\bottomrule
\end{tabular}
\end{table}

For the 10-class task, both pretrained backbones outperformed the scratch CNN on
the test set. EfficientNet-B0 reached a test macro-F1 of 0.9719 and ResNet-50
reached 0.9812, compared with 0.9519 for the scratch CNN. ResNet-50, fine-tuned
through \texttt{layer4}, achieved the best 10-class result overall: it improved
test accuracy by 2.9 percentage points over the scratch CNN (0.9817 vs.\ 0.9531),
with EfficientNet-B0 following closely. This confirms that the pretrained features
transfer well to EuroSAT once the backbones are fine-tuned for enough epochs.

For the binary task, both pretrained models performed better than the scratch
CNN. EfficientNet-B0 reached a test macro-F1 of 0.9941 and ResNet-50 reached
0.9938, compared with 0.9795 for the scratch CNN. The two backbones are
essentially tied on this task, with EfficientNet-B0 marginally ahead. This
suggests that pretrained visual features are especially useful when the target
decision boundary is less fine-grained. Separating vegetation from non-vegetation does not require
distinguishing all visually similar crop and vegetation subclasses.

The data efficiency comparison in Tables \ref{tab:data_efficiency_exp} and
\ref{tab:data_efficiency_binary} is more informative than the full-data numbers
alone. At 500 training images, the ResNet (layer4) achieves a test accuracy of
92.8\% on the 10-class task, while the scratch CNN reaches only 70.2\%. This
22.6 percentage point gap shrinks to 3.4 points at full data. The result shows
that transfer learning is most beneficial in low-data regimes, where the
pretrained features compensate for the lack of labeled examples. When sufficient
data is available, a well-regularized scratch CNN narrows the gap substantially.

The binary data efficiency results are even more striking. At 500 images, the
ResNet already achieves 97.5\% test accuracy, while the scratch CNN reaches
86.9\%. The scratch CNN only surpasses 97\% when trained on the full dataset.
Transfer learning therefore reaches the scratch CNN's full-data performance using
roughly 2.6\% of the training data.

\begin{table}[H]
\centering
\caption{Summary of key comparison metrics between the scratch CNN and ResNet-50.}
\label{tab:comparison_summary}
\small
\begin{tabular}{lrr}
\toprule
Metric & Scratch CNN & ResNet-50 (layer4) \\
\midrule
Trainable parameters          & 390{,}986   & 14{,}985{,}226 \\
Input resolution              & 64$\times$64    & 224$\times$224     \\
Epochs to converge (full data)& 42          & 15                 \\
Time per epoch (full data)    & $\sim$8\,s  & $\sim$151\,s       \\
Total training time           & $\sim$6\,min & $\sim$38\,min     \\
10-class test accuracy        & 0.9531      & 0.9817             \\
10-class test macro-F1        & 0.9519      & 0.9812             \\
Binary test accuracy          & 0.9795      & 0.9938             \\
Test acc.\ at 500 train images (10-class) & 0.7017 & 0.9284      \\
\bottomrule
\end{tabular}
\end{table}

Overall, both approaches are valid. The scratch CNN is easiest to understand
and directly demonstrates the training process from the course. For Model 2,
both transfer-learning backbones use external pretraining knowledge and a larger
input resolution: EfficientNet-B0 is a strong, parameter-efficient baseline that
matches ResNet-50 on the binary task, while ResNet-50 gives the best 10-class
accuracy and by far the best data efficiency, at the cost of higher per-epoch
compute. The fact that both pretrained backbones outperform the scratch CNN on
both tasks indicates that the benefit comes from transfer learning itself rather
than from a single architecture.

\subsection{Limitations}

There are several limitations. First, only RGB information was used. Sentinel-2
data contains additional spectral bands that could be useful for distinguishing
vegetation, water, and built-up areas. Both pretrained backbones expect
three-channel input, so these additional bands were not used in the transfer
learning models.

Second, all models were evaluated with one fixed seed. Multiple training runs
with different seeds would give a better estimate of variance.
